\scanpage{73}
% archive foliation: 72
% lines 1--2
et $S,J$ sont deux parties discrètes
disjointes de $\widehat X$.
% lines 3--8
On va \deleted{\unsure{à}} passer aux catégories
isotopiques correspondantes, et
en donner des interprétations
en termes de théorie des groupes.
Bien entendu, pour cette description,
on est ramené au cas connexe.
% lines 9--16
On peut \deleted{aussi} noter tt de suite
que pour ce pb, on est ramené
au cas où $J=\varnothing$ [en \deleted{\unsure{remplaçant}}\inserted{\unsure{décrivant}}
dans le cas général comme celui
de la donnée d'un $(\widehat X,I)$
\uline{plus} la donnée d'une partie $J$
de $I$].
% lines 17--20
Mais on \unsure{gardera|garde} le
pb sous \inserted{la} forme générale, en
gardant également la dissymétrie
dans le rôle des $S,J$ -- mais
